\chapter{Computational Geometry}\label{ch:geometry}

Geometry problems are mostly a few robust primitives used with care. The golden rule: \textbf{stay in integers whenever the input is
integral}. Orientation tests, areas, convex hulls and closest pairs need only additions, multiplications and comparisons, and then they
are exact. Floating point enters only with circles, angles and intersection points, and there it needs tolerances and a healthy
suspicion. This chapter builds the primitives, the classic algorithms on top of them (hulls, rotating calipers, closest pairs, sweeps),
and two randomised and incremental gems: half-plane intersection and the minimum enclosing circle.

\begin{chapterintro}{\gold\ \plat\ \ioi}{Chapter~\ref{ch:cpp} (integer overflow) and Chapter~\ref{ch:segtree} (Fenwick trees for
sweeps).}
  \item compute orientations, segment intersections, areas and lattice-point counts exactly with cross products;
  \item build convex hulls with the monotone chain, and measure them with rotating calipers;
  \item find the closest pair by a sweep, and handle Manhattan distances by rotating the plane;
  \item design sweep-line algorithms around events and an active set;
  \item intersect half-planes, and find the minimum enclosing circle in expected linear time.
\end{chapterintro}

% ======================================================================================================================
\section[Vectors and the cross product]{Vectors and the cross product \hfill \gold}\label{sec:vectors}

A point is a vector $(x, y)$, and two products of vectors do almost all the work:
\begin{itemize}
  \item the \defn{dot product} $a \cdot b = a_x b_x + a_y b_y = |a|\,|b| \cos\theta$ is positive when the angle between $a$ and $b$ is
        acute, zero when they are perpendicular, negative when it is obtuse;
  \item the \defn{cross product} $a \times b = a_x b_y - a_y b_x = |a|\,|b| \sin\theta$ is positive when $b$ is counter-clockwise from
        $a$, negative when it is clockwise, zero when they are collinear; $|a \times b|$ is the area of the parallelogram they span
        (Figure~\ref{fig:cross}).
\end{itemize}
The \defn{orientation} of three points, $\mathrm{orient}(a, b, c) = \mathrm{sign}\bigl((b - a) \times (c - a)\bigr)$, says whether the
path $a \to b \to c$ turns left ($+1$), right ($-1$), or goes straight (0). It is the single most important primitive of the chapter.

\begin{figure}[htbp]
\centering
\begin{tikzpicture}
  \begin{scope}
    \coordinate (O) at (0, 0);
    \coordinate (A) at (2.4, 0.4);
    \coordinate (B) at (0.9, 1.8);
    \fill[fgsky!25] (O) -- (A) -- ($(A)+(B)$) -- (B) -- cycle;
    \draw[->, line width=1.3pt, draw=fgblue] (O) -- (A) node[right, font=\small] {$a$};
    \draw[->, line width=1.3pt, draw=fgorange!85!black] (O) -- (B) node[above, font=\small] {$b$};
    \node[font=\sffamily\scriptsize, text=fgblue] at (1.65, 1.1) {area $= a \times b > 0$};
  \end{scope}
  \begin{scope}[xshift=6cm]
    \node[vtxs] (p) at (0, 0) {$a$};
    \node[vtxs] (q) at (2.4, 0.5) {$b$};
    \node[vtxs, good] (l) at (1.1, 1.6) {$c$};
    \node[vtxs, hot] (r) at (1.6, -1.2) {$c'$};
    \draw[ded, line width=1pt] (p) -- (q);
    \draw[line width=0.6pt, draw=fggreen, dashed] (q) -- (l);
    \draw[line width=0.6pt, draw=fgverm, dashed] (q) -- (r);
    \node[notel, anchor=west, text width=4.6cm] at (3, 0.2) {$\mathrm{orient}(a, b, c) = +1$: the path turns left to $c$ (green);
      $\mathrm{orient}(a, b, c') = -1$: it turns right (red). Points on the line through $a$ and $b$ give 0.};
  \end{scope}
\end{tikzpicture}
\caption{The cross product is a signed area, and its sign is an orientation.}
\label{fig:cross}
\end{figure}

\usetemplate{Geometry primitives}{Ctrl+Alt+T, Shift+3}{cppoint}

\begin{problem}{Point Location Test (CSES)}
For each of $10^5$ tests, decide whether a point lies to the left of, to the right of, or on the line through two other points
(coordinates up to $10^9$).
\solution One orientation per test: 0.14\,s.
\end{problem}

\begin{code}
void solve() {  // Point, orient (cppoint) are pasted above solve()
    int t;
    cin >> t;
    while (t--) {
        Point a, b, c;
        cin >> a >> b >> c;
        int o = orient(a, b, c);  // the sign of (b - a) x (c - a)
        print(o > 0 ? "LEFT" : o < 0 ? "RIGHT" : "TOUCH");
    }
}
\end{code}

\begin{pitfall}[Precision and overflow]
Coordinates up to $10^9$ give differences up to $2 \cdot 10^9$ and cross products up to $8 \cdot 10^{18}$: they fit in a \cc{long long},
barely. Anything that multiplies three coordinates does not; use \cc{__int128} there. When floating point is unavoidable (circles,
angles, intersection points), use \cc{long double}, compare with a tolerance such as $10^{-9}$ scaled to the magnitudes involved, and
replace \cc{sqrt} and \cc{atan2} by comparisons of squared lengths and cross products wherever possible.
\end{pitfall}

\begin{trick}[Sorting by angle without atan2]
To sort vectors by polar angle exactly, split them into two half-planes, the upper one ($y > 0$, or $y = 0$ and $x > 0$) and the rest;
compare the halves first, and within a half put $a$ before $b$ when $a \times b > 0$. The comparison is integer-exact, it is the basis of
angular sweeps (``the most points inside a $90^\circ$ wedge''), and it also sorts the half-planes of Section~\ref{sec:halfplane}.
\end{trick}

% ======================================================================================================================
\section[Segment intersection]{Segment intersection \hfill \gold}\label{sec:segments}

Segments $ab$ and $cd$ cross properly exactly when $c$ and $d$ lie strictly on opposite sides of the line $ab$ \emph{and} $a$ and $b$ lie
strictly on opposite sides of the line $cd$: $\mathrm{orient}(a, b, c) \cdot \mathrm{orient}(a, b, d) < 0$ and $\mathrm{orient}(c, d, a)
\cdot \mathrm{orient}(c, d, b) < 0$. Every other way of sharing a point puts an endpoint of one segment on the other segment, which is
checked directly; collinear segments that do not overlap must be rejected even though all four orientations are 0
(Figure~\ref{fig:seginter}).

\begin{figure}[htbp]
\centering
\begin{tikzpicture}
  \tikzset{sg/.style={line width=1.4pt, line cap=round}, sp/.style={circle, fill=fgink, inner sep=1.3pt}}
  \foreach \px/\ttx in {0/{proper crossing}, 3.4/{an endpoint on a segment}, 6.8/{collinear, overlapping}, 10.2/{collinear, apart}}
    \node[lbl] at (\px + 1.1, -1.25) {\ttx};
  \draw[sg, draw=fgblue] (0, 0) -- (2.2, 1.0);
  \draw[sg, draw=fgorange!85!black] (0.3, 1.1) -- (1.9, -0.2);
  \draw[sg, draw=fgblue] (3.4, 0) -- (5.6, 0.9);
  \draw[sg, draw=fgorange!85!black] (4.5, 0.45) -- (4.9, -0.5);
  \node[sp] at (4.5, 0.45) {};
  \draw[sg, draw=fgblue] (6.8, 0) -- (8.6, 0.72);
  \draw[sg, draw=fgorange!85!black, dash pattern=on 6pt off 2pt] (7.6, 0.32) -- (9.0, 0.88);
  \draw[sg, draw=fgblue] (10.2, 0) -- (11.0, 0.32);
  \draw[sg, draw=fgorange!85!black] (11.5, 0.52) -- (12.4, 0.88);
  \node[font=\sffamily\scriptsize\bfseries, text=fggreen] at (1.1, 1.45) {YES};
  \node[font=\sffamily\scriptsize\bfseries, text=fggreen] at (4.5, 1.45) {YES};
  \node[font=\sffamily\scriptsize\bfseries, text=fggreen] at (7.9, 1.45) {YES};
  \node[font=\sffamily\scriptsize\bfseries, text=fgverm] at (11.3, 1.45) {NO};
\end{tikzpicture}
\caption{The cases of segment intersection. Only the first is decided by the signs alone; the others need ``does an endpoint lie on the
other segment'', and the last shows why collinearity alone is not enough.}
\label{fig:seginter}
\end{figure}

\usetemplate{Segment intersection}{Ctrl+Alt+Y, 7}{cpseginter}

\begin{problem}{Line Segment Intersection (CSES)}
For each of $10^5$ pairs of segments with integer coordinates up to $10^9$, decide whether they share a point.
\solution Four orientations and four ``point on segment'' tests per pair, all exact: 0.17\,s.
\end{problem}

\begin{code}
void solve() {  // Point (cppoint), segments_intersect (cpseginter) pasted above
    int t;
    cin >> t;
    while (t--) {
        Point a, b, c, d;
        cin >> a >> b >> c >> d;
        yn(segments_intersect(a, b, c, d));
    }
}
\end{code}

When the intersection \emph{point} is needed, solve $a + t(b - a) = c + u(d - c)$: with $r = b - a$ and $s = d - c$, $t = \frac{(c - a)
\times s}{r \times s}$. The point has rational coordinates, so keep the numerator and the denominator as integers if it must be compared
exactly, and only convert to floating point for printing.

% ======================================================================================================================
\section[Polygons]{Polygons \hfill \gold}\label{sec:polygons}

\paragraph{Area.} For a polygon with vertices $p_0, \dots, p_{n-1}$ in order, the \textbf{shoelace formula} gives twice its signed area:
\[
  2A = \sum_{i=0}^{n-1} p_i \times p_{i+1} \qquad (p_n = p_0),
\]
positive for counter-clockwise order. Each term is twice the signed area of the triangle formed by the origin and one edge; triangles
of edges seen counter-clockwise from the origin count positively and the others negatively, and everything outside the polygon cancels
(Figure~\ref{fig:shoelace}).

\begin{figure}[htbp]
\centering
\begin{tikzpicture}[scale=0.62]
  \coordinate (O) at (0, 0);
  \coordinate (A) at (2, 1);
  \coordinate (B) at (6, 1.5);
  \coordinate (C) at (6.5, 4);
  \coordinate (D) at (3, 4.5);
  \fill[fgsky!35, opacity=0.8] (O) -- (A) -- (B) -- cycle;
  \fill[fgsky!35, opacity=0.8] (O) -- (B) -- (C) -- cycle;
  \fill[fgsky!35, opacity=0.8] (O) -- (C) -- (D) -- cycle;
  \fill[fgverm!30, opacity=0.8] (O) -- (D) -- (A) -- cycle;
  \draw[line width=0.4pt, draw=fgline] (O) -- (A) (O) -- (B) (O) -- (C) (O) -- (D);
  \draw[line width=1.4pt, draw=fgink] (A) -- (B) -- (C) -- (D) -- cycle;
  \fill[fgink] (O) circle (2pt) node[below left, lbl] {origin};
  \foreach \pn in {A, B, C, D} \fill[fgink] (\pn) circle (2pt);
  \node[font=\small, below] at (A) {$A$};
  \node[font=\small, below right] at (B) {$B$};
  \node[font=\small, right] at (C) {$C$};
  \node[font=\small, above] at (D) {$D$};
  \node[notel, anchor=west, text width=6.4cm] at (7.6, 2.3) {The polygon $ABCD$ (counter-clockwise). The triangles $OAB$, $OBC$ and
    $OCD$ (blue) have positive signed area, $ODA$ (red) negative; the red one exactly cancels the part of the blue ones that lies
    outside the polygon.};
\end{tikzpicture}
\caption{The shoelace formula as a sum of signed triangles.}
\label{fig:shoelace}
\end{figure}

\usetemplate{Polygon: area, point inside, lattice points}{Ctrl+Alt+Y, 8}{cppolygon}

\begin{problem}{Polygon Area (CSES)}
Print twice the area of a polygon given by its vertices, with coordinates up to $10^9$.
\solution The absolute value of the shoelace sum: twice the area is an integer, so print it without dividing.
\end{problem}

\begin{code}
void solve() {  // Point (cppoint), area2 (cppolygon) are pasted above solve()
    int n;
    cin >> n;
    vector<Point> poly(n);
    for (auto& p : poly) cin >> p;
    print(llabs(area2(poly)));  // the shoelace sum is twice the signed area
}
\end{code}

\paragraph{Point in polygon.} Cast a ray from the point to the right and count the edges it crosses: an odd count means inside
(Figure~\ref{fig:raycast}). Check ``on the boundary'' first, and count an edge only when the point's $y$ lies in the half-open range
$[\min(y_1, y_2), \max(y_1, y_2))$ of the edge, so that a ray through a vertex is counted once, and a horizontal edge never.

\begin{figure}[htbp]
\centering
\begin{tikzpicture}[scale=0.55]
  \fill[fgfill] (0, 0) -- (6, 0) -- (6, 4) -- (4, 4) -- (4, 2) -- (2, 2) -- (2, 4) -- (0, 4) -- cycle;
  \draw[line width=1.2pt, draw=fgink] (0, 0) -- (6, 0) -- (6, 4) -- (4, 4) -- (4, 2) -- (2, 2) -- (2, 4) -- (0, 4) -- cycle;
  \draw[->, line width=0.9pt, draw=fggreen] (1, 3) -- (8, 3);
  \draw[->, line width=0.9pt, draw=fgverm, dashed] (3, 3.3) -- (8, 3.3);
  \fill[fggreen] (1, 3) circle (4pt);
  \fill[fgverm] (3, 3.3) circle (4pt);
  \foreach \cx/\cy in {2/3, 4/3, 6/3} \draw[line width=0.8pt, draw=fggreen] (\cx - 0.2, \cy - 0.2) -- (\cx + 0.2, \cy + 0.2);
  \foreach \cx/\cy in {4/3.3, 6/3.3} \draw[line width=0.8pt, draw=fgverm] (\cx - 0.2, \cy + 0.4) -- (\cx + 0.2, \cy);
  \node[notel, anchor=west, text width=6.4cm] at (8.6, 2.4) {The green point is inside: its ray crosses three edges. The red point,
    in the notch, is outside: its ray crosses two. Only the parity matters, whatever the shape.};
\end{tikzpicture}
\caption{Ray casting: an odd number of crossings means inside.}
\label{fig:raycast}
\end{figure}

\begin{problem}{Point in Polygon (CSES)}
For $m \le 1000$ points and a polygon with $n \le 1000$ vertices, print INSIDE, OUTSIDE or BOUNDARY for every point.
\solution Ray casting in $\Oh{n}$ per point.
\end{problem}

\begin{code}
void solve() {  // Point (cppoint), in_polygon (cppolygon) pasted above
    int n, m;
    cin >> n >> m;
    vector<Point> poly(n);
    for (auto& p : poly) cin >> p;
    while (m--) {
        Point q;
        cin >> q;
        int w = in_polygon(poly, q);  // parity of ray crossings
        print(w > 0 ? "INSIDE" : w < 0 ? "OUTSIDE" : "BOUNDARY");
    }
}
\end{code}

\paragraph{Lattice points.} \textbf{Pick's theorem}: a simple polygon whose vertices are lattice points has area $A = I + \frac{B}{2} - 1$,
where $I$ counts the lattice points strictly inside and $B$ those on the boundary (Figure~\ref{fig:pick}). An edge from $p$ to $q$ contains
$\gcd(|q_x - p_x|, |q_y - p_y|)$ lattice points, counting $q$ but not $p$, so $B$ is a sum of gcds, and $I = (2A - B + 2) / 2$ follows
without enumerating any point.

\begin{figure}[htbp]
\centering
\begin{tikzpicture}[scale=0.6]
  \fill[fgsky!18] (0, 0) -- (6, 0) -- (6, 2) -- (3, 5) -- (0, 3) -- cycle;
  \draw[line width=1.2pt, draw=fgblue] (0, 0) -- (6, 0) -- (6, 2) -- (3, 5) -- (0, 3) -- cycle;
  \foreach \xx in {0, ..., 6} \foreach \yy in {0, ..., 5} \fill[fgline] (\xx, \yy) circle (1.4pt);
  \foreach \xx/\yy in {1/1, 1/2, 1/3, 2/1, 2/2, 2/3, 2/4, 3/1, 3/2, 3/3, 3/4, 4/1, 4/2, 4/3, 5/1, 5/2}
    \fill[fggreen] (\xx, \yy) circle (3.2pt);
  \foreach \xx/\yy in {0/0, 0/1, 0/2, 0/3, 1/0, 2/0, 3/0, 3/5, 4/0, 4/4, 5/0, 5/3, 6/0, 6/1, 6/2}
    \fill[fgverm] (\xx, \yy) circle (3.2pt);
  \node[notel, anchor=west, text width=6.6cm] at (7.3, 2.5) {The pentagon $(0,0), (6,0), (6,2), (3,5), (0,3)$ has $2A = 45$. Its
    boundary carries $B = 15$ lattice points (red), so $I = (45 - 15 + 2)/2 = 16$ lie inside (green): $A = 16 + 7.5 - 1 = 22.5$.};
\end{tikzpicture}
\caption{Pick's theorem counts lattice points from the area and the boundary.}
\label{fig:pick}
\end{figure}

\begin{problem}{Polygon Lattice Points (CSES)}
Count the lattice points strictly inside and on the boundary of a polygon with $n \le 10^5$ vertices and coordinates up to $10^9$.
\solution Pick's theorem with the shoelace area and a gcd per edge: 0.09\,s.
\end{problem}

\begin{code}
void solve() {  // Point, area2, boundary_points (cppoint, cppolygon) pasted
    int n;
    cin >> n;
    vector<Point> poly(n);
    for (auto& p : poly) cin >> p;
    ll twice = llabs(area2(poly)), b = boundary_points(poly);
    print((twice - b + 2) / 2, b);  // Pick's theorem: A = I + B / 2 - 1
}
\end{code}

% ======================================================================================================================
\section[Convex hull]{Convex hull \hfill \plat}\label{sec:hull}

The \defn{convex hull} of a point set is the smallest convex polygon that contains it: a rubber band stretched around the points.
Andrew's \textbf{monotone chain} sorts the points by $x$ (then $y$) and builds the lower hull from left to right: append each point, but
first pop the last point of the chain while the last two chain points and the new one do not turn left. The upper hull is the same
scan from right to left, and the two chains together form the hull (Figure~\ref{fig:hull}). A popped point lies inside the triangle formed
by its neighbours and a point below it, so it can never return, and every point is pushed and popped at most once per chain: $\Oh{n}$
after the $\Oh{n \log n}$ sort.

\begin{figure}[htbp]
\centering
\begin{tikzpicture}[scale=0.85]
  \foreach \px/\py in {3/3, 2/2, 4/3, 1/2, 4/1}
    \fill[fgline] (\px, \py) circle (2.6pt);
  \draw[line width=1.5pt, draw=fgblue] (0, 1) -- (2, 0) -- (6, 2);
  \draw[line width=1.5pt, draw=fgorange!85!black] (6, 2) -- (5, 4) -- (1, 4) -- (0, 1);
  \foreach \px/\py in {0/1, 2/0, 6/2, 5/4, 1/4} \fill[fgink] (\px, \py) circle (3pt);
  \draw[line width=0.7pt, draw=fgverm, dashed] (2, 0) -- (4, 1) -- (6, 2);
  \node[font=\sffamily\scriptsize, text=fgverm, anchor=north west] at (4.05, 0.95) {collinear};
  \node[notel, anchor=west, text width=6.4cm] at (7, 2) {Lower hull (blue) from left to right, upper hull (orange) from right to
    left. The point $(4, 1)$ lies on the segment from $(2, 0)$ to $(6, 2)$: the template pops it, which gives the hull's vertices;
    popping only on strict right turns keeps it, which lists every point on the boundary.};
\end{tikzpicture}
\caption{The monotone chain on ten points: lower and upper hull.}
\label{fig:hull}
\end{figure}

\usetemplate{Convex hull}{Ctrl+Alt+T, Shift+V}{cphull}

\begin{problem}{Convex Hull (CSES)}
Print all of the $n \le 2 \cdot 10^5$ points that lie on the boundary of the convex hull, collinear points included.
\solution The monotone chain, popping only on strict right turns so that collinear points stay. If all the points are collinear, both
chains list all of them; sorting and removing duplicates at the end handles that case and the two shared endpoints. 0.13\,s.
\end{problem}

\begin{code}
void solve() {
    int n;
    cin >> n;
    vector<pll> p(n);
    for (auto& [x, y] : p) cin >> x >> y;
    sort(all(p));
    auto cross = [](pll o, pll a, pll b) {  // > 0: o -> a -> b turns left
        ll ax = a.first - o.first, ay = a.second - o.second;
        return ax * (b.second - o.second) - ay * (b.first - o.first);
    };
    vector<pll> h;
    rep(pass, 2) {  // lower hull left to right, then upper hull back
        int b0 = sz(h);  // where this chain starts
        for (auto q : p) {
            while (sz(h) >= b0 + 2 && cross(h[sz(h) - 2], h.back(), q) < 0)
                h.pop_back();  // pop only right turns: collinear points stay
            h.push_back(q);
        }
        reverse(all(p));
    }
    sort(all(h));  // the chains share their ends; collinear inputs repeat too
    h.erase(unique(all(h)), h.end());
    print(sz(h));
    for (auto [x, y] : h) print(x, y);
}
\end{code}

\begin{trick}[What to do with a hull]
\begin{itemize}
  \item A linear function $ax + by$ is maximised at a hull vertex; with the hull's edge directions sorted, a binary search finds it in
        $\Oh{\log h}$ per query. The convex hull trick of Section~\ref{sec:cht} is this idea for lines.
  \item A point is inside a convex polygon iff it lies in the wedge at vertex 0 found by binary search, on the inner side of that wedge's
        edge: $\Oh{\log n}$.
  \item The farthest pair, the width, and the smallest enclosing rectangle are found by rotating calipers (next section).
\end{itemize}
\end{trick}

% ======================================================================================================================
\section[Rotating calipers]{Rotating calipers \hfill \ioi}\label{sec:calipers}

The two farthest points of a set are both vertices of its convex hull, and they are \emph{antipodal}: some pair of parallel lines through
them encloses the whole hull. \textbf{Rotating calipers} enumerate all antipodal pairs in $\Oh{h}$: for every edge of the hull, in order,
advance a second pointer to the vertex farthest from that edge's line (the area of the triangle grows and then shrinks, so the pointer
only moves forward), and test the distances from the edge's endpoints to that vertex (Figure~\ref{fig:calipers}). The same walk gives the
width of the polygon, the smallest distance between two parallel lines enclosing it, which always has one line flush with an edge.

\begin{figure}[htbp]
\centering
\begin{tikzpicture}[scale=0.85]
  \fill[fgfill] (0, 1) -- (2, 0) -- (6, 2) -- (5, 4) -- (1, 4) -- cycle;
  \draw[line width=1.2pt, draw=fgink] (0, 1) -- (2, 0) -- (6, 2) -- (5, 4) -- (1, 4) -- cycle;
  \draw[line width=1pt, draw=fgblue] ($(0, 1)!-0.7!(2, 0)$) -- ($(0, 1)!1.9!(2, 0)$);
  \draw[line width=1pt, draw=fgblue] ($(5, 4)+(-2, 1)$) -- ($(5, 4)+(2.2, -1.1)$);
  \draw[line width=0.8pt, draw=fgblue, dashed] (0, 1) -- (5, 4) (2, 0) -- (5, 4);
  \draw[line width=1.6pt, draw=fgverm] (0, 1) -- (6, 2);
  \foreach \px/\py in {0/1, 2/0, 6/2, 5/4, 1/4} \fill[fgink] (\px, \py) circle (2.6pt);
  \node[font=\sffamily\scriptsize, text=fgverm, above] at (3.2, 1.55) {diameter $\sqrt{37}$};
  \node[notel, anchor=west, text width=5.8cm] at (7.6, 2) {For the edge from $(0, 1)$ to $(2, 0)$ the farthest vertex is $(5, 4)$, and
    the two blue lines are calipers enclosing the hull. Testing the pairs (dashed) for every edge finds the farthest pair, $(0, 1)$
    and $(6, 2)$ (red).};
\end{tikzpicture}
\caption{Rotating calipers on a convex pentagon.}
\label{fig:calipers}
\end{figure}

\usetemplate{Rotating calipers}{Ctrl+Alt+U, 5}{cpcalipers}

\begin{problem}{Farthest pair}
Find the largest squared distance between two of $n \le 2 \cdot 10^5$ points with coordinates up to $10^9$.
\solution The convex hull, then rotating calipers over its vertices: 0.13\,s, against $\Oh{n^2}$ for all pairs.
\end{problem}

\begin{code}
void solve() {  // convex_hull (cphull), hull_diameter2 (cpcalipers) pasted
    int n;
    cin >> n;
    vector<pll> p(n);
    for (auto& [x, y] : p) cin >> x >> y;
    print(hull_diameter2(convex_hull(p)));  // both ends are hull vertices
}
\end{code}

% ======================================================================================================================
\section[The closest pair of points]{The closest pair of points \hfill \plat}\label{sec:closest}

Sort the points by $x$ and sweep from left to right, keeping the best squared distance $d^2$ found so far and, in a set ordered by $y$,
the points whose $x$ is within $d$ of the sweep line. A new point $p$ can only improve the answer with points of that set whose $y$ lies in
$[p_y - d, p_y + d]$: a $d \times 2d$ box (Figure~\ref{fig:closest}). The points in the box are pairwise at least $d$ apart, so at most 8
of them fit, and each step costs $\Oh{\log n}$: the whole sweep is $\Oh{n \log n}$, with exact integer distances.

\begin{figure}[htbp]
\centering
\begin{tikzpicture}[scale=0.6]
  \fill[fgsky!18] (3.5, 0) rectangle (6, 6);
  \fill[fgorange!25] (3.5, 0.75) rectangle (6, 5.75);
  \draw[line width=0.8pt, draw=fgorange!85!black] (3.5, 0.75) rectangle (6, 5.75);
  \draw[line width=1pt, draw=fgverm, dashed] (6, -0.3) -- (6, 6.3);
  \foreach \px/\py in {0.5/1.2, 1.2/4.6, 2.0/2.5, 2.8/5.6, 3.0/0.6, 4.1/2.2, 4.6/4.4, 5.2/0.4, 5.5/3.9}
    \fill[fgline!80!fgink] (\px, \py) circle (3pt);
  \fill[fgverm] (6, 3.25) circle (4pt) node[right=2pt, font=\sffamily\scriptsize, text=fgverm] {$p$};
  \draw[<->, line width=0.7pt, draw=fgink!70] (3.5, -0.5) -- node[below, lbl] {$d$} (6, -0.5);
  \draw[<->, line width=0.7pt, draw=fgink!70] (7.3, 0.75) -- node[right, lbl] {$2d$} (7.3, 5.75);
  \node[notel, anchor=west, text width=6.4cm] at (8.6, 3) {Only the points left of the sweep line by less than $d$ (blue strip) are
    kept, and of those only the ones in the orange box can be closer to $p$ than $d$. Points at mutual distance at least $d$ fit
    into a $d \times 2d$ box at most eight at a time.};
\end{tikzpicture}
\caption{The closest-pair sweep checks a constant number of candidates per point.}
\label{fig:closest}
\end{figure}

\usetemplate{Closest pair}{Ctrl+Alt+Y, 9}{cpclosest}

\begin{problem}{Minimum Euclidean Distance (CSES)}
Print the smallest squared distance between two of $n \le 2 \cdot 10^5$ points.
\solution The sweep, in exact integers: 0.13\,s.
\end{problem}

\begin{code}
void solve() {  // Point (cppoint), closest_pair (cpclosest) pasted above
    int n;
    cin >> n;
    vector<Point> p(n);
    for (auto& q : p) cin >> q;
    print(closest_pair(p));  // squared, so it stays exact
}
\end{code}

% ======================================================================================================================
\section[Manhattan distances]{Manhattan distances \hfill \gold}\label{sec:manhattan}

The Manhattan distance $|x_1 - x_2| + |y_1 - y_2|$ becomes much easier after rotating the plane by $45^\circ$. With $u = x + y$ and $v
= x - y$,
\[
  |x_1 - x_2| + |y_1 - y_2| = \max(|u_1 - u_2|, |v_1 - v_2|),
\]
because $|a| + |b| = \max(|a + b|, |a - b|)$. The diamond of points within Manhattan distance $r$ becomes an axis-parallel square
(Figure~\ref{fig:manhattan}), and the two coordinates separate: the farthest pair, nearest-neighbour queries and ``all points within
distance $r$'' become independent questions about $u$ and $v$.

\begin{figure}[htbp]
\centering
\begin{tikzpicture}[scale=0.55]
  \begin{scope}
    \draw[->, line width=0.5pt, draw=fgink!60] (-3, 0) -- (3.2, 0) node[right, lbl] {$x$};
    \draw[->, line width=0.5pt, draw=fgink!60] (0, -3) -- (0, 3.2) node[above, lbl] {$y$};
    \fill[fgsky!30] (2, 0) -- (0, 2) -- (-2, 0) -- (0, -2) -- cycle;
    \draw[line width=1.2pt, draw=fgblue] (2, 0) -- (0, 2) -- (-2, 0) -- (0, -2) -- cycle;
    \node[lbl] at (0, -3.6) {$|x| + |y| \le 2$};
  \end{scope}
  \draw[->, line width=1pt, draw=fgink!60] (4.5, 0) -- node[above, lbl] {$u = x + y$, $v = x - y$} (8.4, 0);
  \begin{scope}[xshift=12.4cm]
    \draw[->, line width=0.5pt, draw=fgink!60] (-3, 0) -- (3.2, 0) node[right, lbl] {$u$};
    \draw[->, line width=0.5pt, draw=fgink!60] (0, -3) -- (0, 3.2) node[above, lbl] {$v$};
    \fill[fgsky!30] (-2, -2) rectangle (2, 2);
    \draw[line width=1.2pt, draw=fgblue] (-2, -2) rectangle (2, 2);
    \node[lbl] at (0, -3.6) {$\max(|u|, |v|) \le 2$};
  \end{scope}
\end{tikzpicture}
\caption{Rotating by $45^\circ$ turns Manhattan balls into squares.}
\label{fig:manhattan}
\end{figure}

\begin{problem}{Maximum Manhattan Distances (CSES)}
Points arrive one at a time ($n \le 2 \cdot 10^5$); after each, print the largest Manhattan distance between two points so far.
\solution In rotated coordinates the answer is $\max(u_{\max} - u_{\min}, v_{\max} - v_{\min})$: four running extremes. 0.14\,s.
\end{problem}

\begin{code}
void solve() {
    int n;
    cin >> n;
    ll umin = INF, umax = -INF, vmin = INF, vmax = -INF;
    rep(i, n) {
        ll x, y;
        cin >> x >> y;
        ll u = x + y, v = x - y;  // |dx| + |dy| = max(|du|, |dv|)
        chmin(umin, u), chmax(umax, u), chmin(vmin, v), chmax(vmax, v);
        print(max(umax - umin, vmax - vmin));
    }
}
\end{code}

\paragraph{Sums of distances.} The \emph{sum} of Manhattan distances over all pairs separates by axis without any rotation: sort the
$x$-coordinates, and the $i$-th smallest contributes $x_i$ for each of the $i$ smaller ones minus their sum, the contribution technique of
Chapter~\ref{ch:patterns}.

\begin{problem}{All Manhattan Distances (CSES)}
Sum the Manhattan distances between all pairs of $n \le 2 \cdot 10^5$ points with coordinates up to $10^9$ in absolute value.
\solution Sort each axis and add up the contributions. The total reaches $2.7 \cdot 10^{19}$ on random inputs at the limits, more than a
\cc{long long} holds, so the sum is kept in \cc{__int128} and printed by hand. 0.15\,s.
\end{problem}

\begin{code}
void solve() {
    int n;
    cin >> n;
    vi xs(n), ys(n);
    rep(i, n) cin >> xs[i] >> ys[i];
    __int128 total = 0;  // can exceed 64 bits: up to about 4e19
    for (vi* c : {&xs, &ys}) {  // |dx| + |dy|: the two axes add up separately
        sort(all(*c));
        __int128 pre = 0;  // the sum of the smaller coordinates
        rep(i, n) total += (__int128)(*c)[i] * i - pre, pre += (*c)[i];
    }
    string s;
    for (; total > 0; total /= 10) s += char('0' + total % 10);
    if (s.empty()) s = "0";
    reverse(all(s));
    print(s);
}
\end{code}

% ======================================================================================================================
\section[Sweep lines]{Sweep lines \hfill \plat}\label{sec:sweep}

A sweep line moves across the plane and processes \emph{events}, such as segment endpoints, points and queries, in order of position,
while a data structure holds the \emph{active} objects that the line currently crosses. It turns a two-dimensional problem into a
sequence of one-dimensional ones. Designing a sweep means answering three questions: what are the events, in which order do events at the
same position happen, and what must the active set support? Section~\ref{sec:mincount} computed the area of a union of rectangles this
way, with a segment tree as the active set.

\begin{figure}[htbp]
\centering
\begin{tikzpicture}[scale=0.55]
  \foreach \xa/\xb/\yy in {0/5/1, 2/9/3, 1/4/5, 6/10/4.5, 3/8/6} \draw[line width=1.4pt, draw=fgblue, line cap=round] (\xa, \yy) -- (\xb, \yy);
  \foreach \xx/\ya/\yb in {1.5/0.5/2, 7/2/5.5, 9.5/0/6.5} \draw[line width=1.4pt, draw=fgorange!85!black, line cap=round] (\xx, \ya) -- (\xx, \yb);
  \draw[line width=1pt, draw=fgverm, dashed] (7, -0.4) -- (7, 7);
  \foreach \yy in {3, 4.5, 6} \fill[fgverm] (7, \yy) circle (3.5pt);
  \node[font=\sffamily\scriptsize, text=fgverm, above] at (7, 7) {sweep at $x = 7$};
  \node[notel, anchor=west, text width=6.0cm] at (11.2, 3.2) {At $x = 7$ the active horizontals (crossing the line) are at heights 3, 4.5
    and 6. The vertical segment from 2 to 5.5 meets the two whose height lies in $[2, 5.5]$: a range count in a Fenwick tree over the
    heights.};
\end{tikzpicture}
\caption{Counting intersections of horizontal and vertical segments by a sweep.}
\label{fig:intpoints}
\end{figure}

\begin{problem}{Intersection Points (CSES)}
$n \le 10^5$ horizontal and vertical segments with coordinates up to $10^6$ are given, and parallel segments never touch. Count the
points where a horizontal and a vertical segment meet.
\solution Sweep over $x$. A horizontal segment becomes active at its left end and inactive at its right end; a vertical segment at $x$
asks how many active horizontals have their height in its $y$-range (Figure~\ref{fig:intpoints}). At equal $x$, insertions come before
queries and queries before removals, so that touching endpoints count. 0.12\,s.
\end{problem}

\begin{code}
void solve() {
    int n;
    cin >> n;
    const ll OFF = 1000001, N = 2 * OFF + 1;  // y + OFF lies in [1, N]
    vector<array<ll, 4>> ev;  // (x, type, y1, y2): 0 add, 1 query, 2 remove
    rep(i, n) {
        ll x1, y1, x2, y2;
        cin >> x1 >> y1 >> x2 >> y2;
        if (y1 == y2) {  // horizontal: alive for x in [min, max]
            ev.push_back({min(x1, x2), 0, y1, y1});
            ev.push_back({max(x1, x2), 2, y1, y1});
        } else ev.push_back({x1, 1, min(y1, y2), max(y1, y2)});
    }
    sort(all(ev));  // at equal x: insert, then query, then remove
    vi bit(N + 1);
    auto add = [&](ll y, ll v) {
        for (ll i = y + OFF; i <= N; i += i & -i) bit[i] += v;
    };
    auto pre = [&](ll y) {  // active horizontals at height <= y
        ll s = 0;
        for (ll i = y + OFF; i > 0; i -= i & -i) s += bit[i];
        return s;
    };
    ll ans = 0;
    for (auto [x, type, a, b] : ev) {
        if (type == 0) add(a, 1);
        else if (type == 2) add(a, -1);
        else ans += pre(b) - pre(a - 1);
    }
    print(ans);
}
\end{code}

\begin{recognize}
Many intervals, rectangles or segments, and a question about their overlaps, unions or intersections; ``for each point, how many
rectangles contain it''; ``the largest number of overlapping intervals'' (a one-dimensional sweep over $+1$ and $-1$ events). Sort the
events, fix the tie-breaking order, and choose the active set's structure: a counter, a multiset, a Fenwick tree, or a segment tree.
\end{recognize}

% ======================================================================================================================
\section[Half-plane intersection]{Half-plane intersection \hfill \ioi}\label{sec:halfplane}

A directed line keeps the half-plane on its left. The intersection of $n$ half-planes is a convex polygon (possibly empty or unbounded),
and it is computed in $\Oh{n \log n}$ much like a convex hull: sort the lines by direction angle, then add them in that order to a deque.
A new line makes the last line in the deque redundant when the intersection point of the last two lines lies outside the new half-plane,
so pop from the back while that happens; the same check at the front removes lines that the new one cuts off from the other end. Of two
lines with the same direction only the stricter one is kept. A final cleanup between the two ends of the deque closes the polygon
(Figure~\ref{fig:halfplane}). Adding a bounding box of four lines keeps the result bounded.

\begin{figure}[htbp]
\centering
\begin{tikzpicture}[scale=0.65]
  \fill[fggreen!25] (1, 1) -- (5, 0.5) -- (6.2, 3.2) -- (3.8, 5.2) -- (0.8, 3.8) -- cycle;
  \draw[line width=1.4pt, draw=fggreen!80!black] (1, 1) -- (5, 0.5) -- (6.2, 3.2) -- (3.8, 5.2) -- (0.8, 3.8) -- cycle;
  \foreach \ax/\ay/\bx/\by in {1/1/5/0.5, 5/0.5/6.2/3.2, 6.2/3.2/3.8/5.2, 3.8/5.2/0.8/3.8, 0.8/3.8/1/1} {
    \draw[->, line width=0.8pt, draw=fgblue] ($(\ax, \ay)!-0.35!(\bx, \by)$) -- ($(\ax, \ay)!1.35!(\bx, \by)$);
  }
  \draw[->, line width=0.8pt, draw=fgline, dashed] (-0.5, 6.2) -- (7, 6.0);
  \node[font=\sffamily\scriptsize, text=fggray, anchor=south] at (6.0, 6.05) {redundant};
  \node[notel, anchor=west, text width=6.0cm] at (8.0, 2.8) {Six half-planes, each to the left of its directed line. Five of them bound
    the intersection (green); the dashed one contains it entirely and is dropped when the deque is cleaned up.};
\end{tikzpicture}
\caption{Half-plane intersection: the lines sorted by angle form the boundary in order.}
\label{fig:halfplane}
\end{figure}

\usetemplate{Half-plane intersection}{Ctrl+Alt+U, 4}{cphalfplane}

\begin{problem}{Half-plane intersection area}
Print the area of the intersection of $n \le 10^5$ half-planes (each to the left of a directed line through two integer points) inside the
square $[-10^6, 10^6]^2$.
\solution The deque algorithm with the square's four sides added, then the shoelace formula. With $10^5$ tangent lines of a circle,
which all survive, it takes 0.83\,s, mostly for the angles and the long-double intersections.
\end{problem}

\begin{code}
void solve() {  // half_plane_intersection (cphalfplane) is pasted above solve()
    int n;
    cin >> n;
    vector<array<ld, 4>> hs(n);
    for (auto& [ax, ay, bx, by] : hs) cin >> ax >> ay >> bx >> by;
    auto poly = half_plane_intersection(hs, 1e6);  // keeps the left of a -> b
    ld area = 0;  // the shoelace formula on the resulting convex polygon
    rep(i, sz(poly)) {
        auto [x1, y1] = poly[i];
        auto [x2, y2] = poly[(i + 1) % sz(poly)];
        area += x1 * y2 - x2 * y1;
    }
    cout << fixed << setprecision(6) << area / 2 << '\n';
}
\end{code}

% ======================================================================================================================
\section[The minimum enclosing circle]{The minimum enclosing circle \hfill \ioi}\label{sec:mincircle}

The smallest circle that contains $n$ given points is unique, and it is determined by two points on its boundary (a diameter) or by
three. \textbf{Welzl's algorithm}, in its iterative form, shuffles the points and maintains the smallest circle of the points seen so far.
When a new point $p_i$ falls outside, the new circle must pass through $p_i$; rebuild it from the earlier points with $p_i$ fixed on the
boundary, by the same procedure one level down, and when a second point $p_j$ falls outside, with both $p_i$ and $p_j$ on the boundary; at
the third level a circle through three points is determined (Figure~\ref{fig:mincircle}).

The rebuilds look expensive, but in random order they are rare. By \emph{backwards analysis}, the circle of the first $i$ points changes
when $p_i$ is added only if $p_i$ is one of the at most three points that define it, which happens with probability at most $3/i$; a
rebuild costs $\Oh{i}$, so its expected cost is $\Oh{1}$ per point, and the same argument one level down keeps the whole algorithm at
expected $\Oh{n}$.

\begin{figure}[htbp]
\centering
\begin{tikzpicture}[scale=0.7]
  \draw[line width=1.3pt, draw=fgblue, fill=fgsky!15] (3.8043, 3.0217) circle (2.9847);
  \foreach \px/\py in {3/3, 4/3, 2/4, 5/4} \fill[fgline!80!fgink] (\px, \py) circle (3pt);
  \foreach \px/\py in {1/2, 6/1, 4/6} \fill[fgverm] (\px, \py) circle (3.8pt);
  \draw[line width=0.7pt, draw=fgverm, dashed] (1, 2) -- (6, 1) -- (4, 6) -- cycle;
  \node[notel, anchor=west, text width=6.0cm] at (7.5, 3) {Seven points; the smallest enclosing circle passes through the three red ones,
    the circumcircle of their triangle, which is acute. Had the triangle been obtuse, the circle would instead have the longest side as
    a diameter and pass through only two points.};
\end{tikzpicture}
\caption{The minimum enclosing circle is determined by two or three of the points.}
\label{fig:mincircle}
\end{figure}

\usetemplate{Minimum enclosing circle}{Ctrl+Alt+U, Shift+M}{cpmincircle}

\begin{problem}{Minimum enclosing circle}
Print the centre and the radius of the smallest circle containing $n \le 10^5$ points with integer coordinates up to $10^6$.
\solution Welzl's algorithm after a shuffle: 0.51\,s, almost all of it in \cc{hypotl}. Comparing squared distances instead would save most
of that.
\end{problem}

\begin{code}
void solve() {  // min_enclosing_circle (cpmincircle) is pasted above solve()
    int n;
    cin >> n;
    vector<pair<ld, ld>> p(n);
    for (auto& [x, y] : p) cin >> x >> y;
    auto [cx, cy, r] = min_enclosing_circle(p);  // expected O(n), shuffled
    cout << fixed << setprecision(6) << cx << ' ' << cy << ' ' << r << '\n';
}
\end{code}

\begin{pitfall}[Shuffle, or be hacked]
Without the random shuffle, Welzl's algorithm is quadratic on inputs sorted by distance from the centre, and such inputs are easy to
construct. Shuffle with a seeded generator (Section~\ref{sec:random-stl}), never with the default seed.
\end{pitfall}

\begin{summary}[Chapter summary]
\begin{itemize}
  \item The cross product gives signed areas and orientations; stay in integers, and mind that cross products of $10^9$-sized coordinates
        nearly fill a \cc{long long}.
  \item Segment intersection is four orientations plus endpoint-on-segment checks; the shoelace formula gives areas; ray casting decides
        point in polygon; Pick's theorem counts lattice points from the area and the boundary.
  \item The monotone chain builds the convex hull in $\Oh{n \log n}$; rotating calipers find the farthest pair and the width in $\Oh{h}$.
  \item The closest pair falls to a sweep with a $y$-ordered set; rotating by $45^\circ$ turns Manhattan distances into maxima of two
        independent coordinates.
  \item A sweep line needs its events, a tie-breaking order and an active set.
  \item Half-planes intersect in $\Oh{n \log n}$ with a deque over lines sorted by angle; the minimum enclosing circle takes expected
        $\Oh{n}$ after a random shuffle.
\end{itemize}
\end{summary}

\begin{practice}
\ppgroup{Geometry}
\begin{pplist}
  \item Point Location Test
  \item Line Segment Intersection
  \item Polygon Area
  \item Point in Polygon
  \item Polygon Lattice Points
  \item Minimum Euclidean Distance
  \item Convex Hull
  \item Maximum Manhattan Distances
  \item All Manhattan Distances
  \item Intersection Points
  \item Line Segments Trace I
  \item Line Segments Trace II
  \item Lines and Queries I
  \item Lines and Queries II
  \item Area of Rectangles
  \item Robot Path
\end{pplist}
\end{practice}
